\chapter*{跋}\addcontentsline{toc}{chapter}{跋}\markboth{跋}{跋}
% 原书第345页

2025年5月5日晚，美国ABC电视台做了一档节目，介绍两天后将开始的新教皇选举。\NoteRef{08}{1}来自世界各地的一百三十三位红衣主教将要在梵蒂冈圣彼得大教堂旁边的西斯廷小教堂中投票选出新教皇。获得三分之二绝对多数选票的主教将成为新教皇。所有进去投票的主教都会被没收手机等通讯设备。如果投票结果没有达到绝对多数，他们将会被一直关在教堂中，吃住都在里面，一轮又一轮地投票，直到选出新教皇。


% 原书第345–346页

近千年来，天主教教皇的每一次选举过程都有详细记录，为ABC台这档节目提供了大量历史素材，让这个看似没有什么特别的选举增添了不少戏剧性。历史上有记载最长的一次选举持续了近三年（1268年12月1日至1271年9月1日），也就是说，各方主教在新教皇人选上长达三年都不能达成一致。从1274年开始，为了加快选举速度，如果三天都没有选出新教皇，那么关在教堂里面的每位主教的中晚餐都将只有一道菜可吃。假设如此吃了五天后还没有选出，那就只有面包、红酒和水提供了。\NoteRef{08}{2}即便如此，19世纪以来还是有多次选举持续了数周。2025年这次选举与以往一样，到底多久能选出新教皇，事前谁也不知道。投票结果的公布依然沿用传统方式—由西斯廷小教堂上面的烟囱来对外发布信号。如果烟囱冒出的是黑烟，就表示这轮投票没有选出新教皇；而如果是白烟，就代表新教皇已选出。这样的选举形式从1274年确定下来后，到现在的七百多年时间里，共举行过七十六次。一直到2025年这次选举，食物依然不允许有整鸡和意大利饺（ravioli）出现，因为这些食物在历史上曾经被用来传递纸条。在讲完选举历史中的趣事后，节目主持人让大家关注明天的选举，并告诉大家会直播烟囱冒烟环节，最后说让我们一起期待“历史的上演（history unfolds）”。


% 原书第346–347页

“历史的上演”是英文的直译，在中文中比较接近的说法是“见证历史时刻的到来”。无论在中文还是西方的语境中，“历史”通常联系的是过去。很多时候，“历史”和“过去”两个词可以互相切换使用。但“历史”还可以用来指尚未发生的事情，这时“历史”和“过去”就不能互换了。比如“见证过去时刻的到来”，就是逻辑上完全不成立的句子。当两个词不能切换时，其实就意味着“历史”有一层“过去”不具有的含义，也就是说这里蕴含着二者的区别。


% 原书第347页

时间顺序上讲，有了过去，然后才会在其中产生历史，历史只是过去发生之事中极少的一部分。那历史又是如何诞生和形成的？这是西方历史哲学中的重要问题。20世纪以来的主流观点都认同，历史是从过去中筛选而来。\NoteRef{08}{3}比如在记录方式极端有限的情况下，古人把他们认为值得记录的大事刻在了青铜器上。至于哪些值得记录，则是当时人的主观筛选。后来一代又一代人，又在最初留下的记录中，进一步筛选出自己认为有价值的部分，同时也是自己关注的部分，书写出一个个有内在逻辑的历史叙事。用现代的理论话语来说，我们看到的历史是在一次又一次的选择后被呈现出来的。


% 原书第347页

历史是被选择书写的过去。这是“历史的上演”或者“见证历史时刻的到来”这些说法能够成立背后的逻辑预设。因为当人们用到这种表达的时候，就是在预言将要发生之事必然被选择记录、书写，成为历史。ABC节目的制作团队在见过了丰富的教皇选举历史材料后，自然能够确信2025年这次选举依然会被详细记录，同时未来的人也会在各种不同场景的历史书写中提到这次选举。所以这次教皇选举用“历史的上演”来表达，恰如其分，也符合大众的认知。


% 原书第347–348页

尽管“历史的上演”是一种可以被理解的说法，但它并不是绝大部分人生活中会有的意识。古往今来绝大部分人的一生都只会成为过去，而不会有哪怕生命的瞬间能成为历史的一部分，因而大家在日常生活中也不会思考被后世选择、成为历史一部分的可能。但是这一点恰好是皇帝以及传统精英在乎的部分，并且可能影响到他们的所言所行。在西方的宗教传统中，肉身终将消逝成为过去，而死后的审判（judgment）才是永恒。孔子说：“知我罪我，其惟春秋。”这里的“春秋”是历史记录的代称。这句话的意思是说，真正理解他的是后世回看过去时做出的历史评判。\NoteRef{08}{4}这其中还暗含一层意思，就是当世的评论与后世的评论可能完全不一样，而后世的评论有可能才是永恒的。屈原的“余将董道而不豫兮，固将重昏而终身”，是在书写他在坚持着不被当世理解的信念。在这句诗之前，他提到了伍子胥、比干这些历史人物，认为不被理解只是“阴阳易位，时不当兮”，是暂时的，同时相信未来的人在回看时，最终会理解他。


% 原书第348页

在中国传统文化中，历史是一个独立于现世以外的永恒空间。肉身不过百年，即便今世不能被理解，到了未来某个时候也有可能被理解。这其实类似于一种信仰。就是贵为皇帝，号令天下的康熙，他也觉得自己有些措施并不被理解，说出了这样的话：


% 原书第348–349页

\begin{bookquote}孔子，圣人也。犹言“知我者，其惟春秋乎；罪我者，其惟春秋乎”。孟子又言：“尽信书则不如无书。”当今之世，用人行政、规模法度之是非，朕当自任，无容他诿。\NoteRef{08}{5}\end{bookquote}

% 原书第349页

康熙相信，是非对错后世自有定论，他也不急于一时之理解。近代历史学家陈寅恪在艰难环境下写出了他的最后一本书《柳如是别传》，无论内容还是理论都不被理解。在书的末尾，留给后世最后的文字中，他“合掌说偈”，最后五句是：


% 原书第349页

\begin{bookquote}
繁琐冗长，见笑君子。\\
失明膑足，尚未聋哑。\\
得成此书，乃天所假。\\
卧榻沉思，然脂瞑写。\\
痛哭古人，留赠来者。\NoteRef{08}{6}
\end{bookquote}

% 原书第349–350页

“偈”是佛门中的一种文体，类似于诗歌。禅门公案中说偈的都是得道高僧，用偈来表达得证果位后的领悟。陈寅恪一生以诗证史，留下了大量诗作。在其文字写作的最后，用偈而不用诗，这是他自信悟到了历史研究的真味，要像高僧临终“合掌说偈”那样来宣示心迹，传承法脉。他清楚《柳如是别传》“繁琐冗长”，也知道当时动荡年代流行的历史研究理论和套路，所以他告知世间，尽管他看不见了，也走不动了，但“尚未聋哑”，他明白新作会“见笑”于时人君子。但“得成此书，乃天所假”，能成此书，是上天眷顾，是上天假借他的手写出。敢于把自己的作品比作天意，是写作之人最高的自我评价。“卧榻沉思，然脂瞑写”一句，表明这本书是其苦心孤诣之作，后世不用推敲这是不是他深思后的作品。写到这里，他写下了“痛哭古人，留赠来者”八个字，直接跳过了今人。也就是说，他相信后世自有历史公论。本质上，陈寅恪与屈原所写并无不同，都是对未来的历史评价有公论抱有信心。


% 原书第350–351页

历史公论存在于后世反思过去而写出的历史记述中。中国系统的历史评论最晚从孔子作《春秋》就开始了。孔子根据鲁国的史料，选择出他认为有意义的部分；然后再用特定的字词对所写事件进行评判。\NoteRef{08}{7}孟子所谓：“孔子成《春秋》而乱臣贼子惧”，就是立足于《春秋》的褒贬评论之上。乱臣贼子并不是惧怕孔子或者史官个人，而是担心他们的事迹被选择性地记述后，成为了历史的一部分，留下永恒的骂名。孔子《春秋》中评判的方法称为“笔削”。“笔”是挑选出书写的部分，“削”是有意识不写的部分。用今天的话来说，孔子就是用选择记叙和回避书写两种方式来针砭历史中的人物和事件。从司马迁说“至于为《春秋》，笔则笔，削则削，子夏之徒不能赞一辞”，到清代章学诚的“史之大原本乎《春秋》，《春秋》之义昭乎笔削”，都一致同意孔子的“笔削”。也就是说，选择性记述是中国传统历史书写的特性。


% 原书第351页

孔子选择性记述后面有一套他要推广的道德标准，即所谓义理。《左传》、《公羊传》和《穀梁传》—俗称的《春秋》三传就是在《春秋》记录事件的框架下，解释孔子选取各个事件，及其措辞背后所反映的义理。比如《春秋》中有“冬，十有二月，祭伯来”一条。《春秋》三传在这一条后的注解基本相同，都着重诠释用“来”字而不用“朝”字背后所蕴含的大义，即祭伯在未得王命允许的前提下，私自出去会见诸侯，他这一不符礼仪的举动就只能称作“来”，不能称作“朝”。孔子选择记录“祭伯来”这件事，以及用哪个字来书写，都是在贯彻他要宣传的道理。孔子在历史的编撰中对过去事件赋予了一层他想要赋予的价值，此即孔子对过去人和事的评判。


% 原书第351–352页

在传统历史书写中，到底是对过去之事的记录重要，还是在事件上赋予的义理重要？ 这个问题很难有一个绝对的答案，因为绝大多数的历史人物和历史事件都只有唯一的记录，我们后世之人无从判断这些记录究竟是准确记录了发生之事，还是为了宣扬某些道理而改变了对事件的描述。就像如果只有一张照片，告诉你这个人就是历史中的张三，那我们没法判断到底这张照片是不是张三。而我们唯一能做的就是接受照片上是张三这个结论。但在一些具体的案例中，因为有了其他的材料，就可以看到义理在历史记录中占有的重要位置。西方中国上古史专家夏含夷（Edward L. Shaughnessy）就曾经写到在对武王伐纣这样重大历史事件的书写上，义理主导了对事件的叙述。这则用义理来指导历史书写的句子正好还是大众耳熟能详的名句：


% 原书第352页

\begin{bookquote}尽信书则不如无书。吾于《武成》，取二三策而已矣。仁者无敌于天下，以至仁伐至不仁，而何其血之流杵也。\end{bookquote}

% 原书第352页

这里孟子否定已有《武成》记录中对周灭商时“血之流杵”惨况的描述，而用儒家宣讲的“仁者无敌于天下”的道理来驳斥。在上千年的时间里，周灭商的主流历史叙述都是在孟子的框架下进行。\NoteRef{08}{8}


% 原书第352–353页

历史叙述的形成并不是写作者单方面的选择。孟子只是讲出了一种他认为正确的历史叙述。但他的版本要成为被后世接受的历史叙述，还需要一代代后来人都认可他的说法才有可能。那制约着历史叙述是否被后世接受，也就是其最后的形成又有哪些因素呢？海登·怀特（Hayden White）给出过一个思路。他提出历史叙述的形成受制于三个方面：原始记录（unprocessed historical record），其他版本的历史叙述（other historical accounts）和读者（audience）。这三个方面共同构成了历史叙述形成的场域（field）。\NoteRef{08}{9}这个场域理论应该是借鉴了市场经济理论。各种历史叙述就是市场中的产品，而读者就是消费者，他们在不停对比历史叙述与原始记录，从中选出他们认同的版本。如果借用海登·怀特的场域理论，那么可以看到孟子在提出他的历史叙述时，也察觉到了上面所说的三点。首先他思考过已有的原始记录《武成》，同时也思考过其他版本的历史叙述，因而他在警告读者不能“尽信书”，也就是不能完全相信书中的叙述。而孟子最初的读者就是他自己的门徒，或者说是最早接受儒家学说的人。


% 原书第353页

孟子提出他的论述时在作比较，当然后世的读者也在比较各家叙述。夏含夷特别提到了司马迁。根据司马迁撰写《史记》用到过的材料，很明显他当时是看到过描述周灭商惨状的文字，但他选择了回避，采信了孟子的说法。司马迁做出了他的选择。而比司马迁晚一些的王充，在比较了各种说法后，就对孟子的叙述逻辑提出了质疑。他认为汉高祖伐秦，后来灭项羽，都是“战场流血，暴尸万数”，因而他认为“周兵不血刃”就把商给灭了的叙述，不符合常理。然后，他说，查阅《武成》，知道了周灭商的“牧野之战，赤地千里”，因而，他认为这种说法才是历史真正发生过的样子。王充比较后，做出了他的判断。\NoteRef{08}{10}


% 原书第353–354页

海登·怀特场域理论成立的关键是读者的自由选择。在这种情况下，历史的形成也像市场经济理论所言，有一只看不见的手，在选择着历史叙述。但现实是在漫漫的历史长河中，历史叙述的流传往往都受到外力的干预，比如西方有宗教意识形态的介入，而中国是孔孟之道独占鳌头。这也就决定了孟子对历史的叙述及其特别筛选出的史料，伴随着儒家思想成为主流的进程而逐渐被广泛接受。与孟子观点不匹配的记录，比如载有周灭商惨状的《武成》，也就是孟子和王充看到过的那个版本，逐渐从古文《尚书》中消失了，被换上了另一篇没有周灭商惨状的文字，却袭用了《武成》的名字。\NoteRef{08}{11}这样后世读到十三经之一的《尚书》时，里面就只有与孟子叙述相匹配的原始资料。孟子的历史叙述也就在官方意识形态的帮助下，成为了后世读者的“选择”。


% 原书第354–355页

除了承认历史叙述是被多重选择后的结果以外，西方历史哲学中还有一个问题，即历史是否可以成为一种知识（knowledge）？这里所谓的知识，是柏拉图提出的超越时间与空间那种形而上的知识。比如科学知识、数学知识就都不受时间、空间的影响。三百年前总结出的公式和实验，现在依然可以重复。但问题是，历史是发生在过去的时间与空间中的事，他们并不能被重复验证，那么历史就肯定受制于时空，它又如何成为知识（knowledge）？18世纪晚期开始，以黑格尔等德国思想家为代表的欧洲学者都在尝试着回答这个问题，试图找出历史中超越时间与空间那部分知识，也就是历史中的规律。在这种思维的引导下，历史被认为与科学相似。当然这个尝试在理论和实际层面都遇到了各种问题，到了19世纪末、20世纪初，以“一切历史都是当代史”而闻名的意大利史学家克罗齐（Benedetto Croce）对历史为什么不是科学给出过一个解释：


% 原书第355页

\begin{bookquote}历史针对的是一个个具体的事实；当然在处理特殊的个体时也会把这些个体联系到背后整体（the whole）中。但这个过程与科学不同。因为历史中的整体性并不是科学追求的普遍性（the general）。\NoteRef{08}{12}\end{bookquote}

% 原书第355–356页

他用强调整体性与普遍性的区别（the whole is different from the general）把历史与科学区分开来。科学是对普遍性规律的总结和揭示，而历史是对过去具体事件的叙述（narrates），是对过去的重建（reproduction）。在这个过程中，最重要的是叙述历史的那个人。从克罗齐以后，西方史学理论开始注重历史写作者的主观性是历史叙述中不可分割的部分。在克罗齐的解释中，历史是写作者根据自己当下的认知而对过去留下的各种材料进行的解读。\NoteRef{08}{13}克罗齐的“一切历史都是当代史”本质上就是强调了所有历史叙述诞生的过程中无法避免的主观性。\NoteRef{08}{14}后来，20世纪西方史学理论也全面放弃了对纯粹历史客观性的执着追求，认为那只是一个曾经有过的“高贵的梦想”。\NoteRef{08}{15}到了后现代主义史学的解释中，由于更多强调历史写作主观性的一面，历史甚至成为了一种文学形式，因为后现代主义学者认为史家的叙述是在历史材料的基础上对过去的想象。\NoteRef{08}{16}总结起来，关于历史的西方理论中，无论是之前强调客观性，将其类比于科学；还是后来强调主观性，将其类比于艺术或者文学，历史在西方都是被当作了被研究的客体。


% 原书第356页

而在中国，历史是社会运作的一部分。历史从来就不是客体，它始终与各个时代的社会运作紧密相关。这是西方史学理论难以涵盖到中国历史书写的原因。20世纪西方总结前人历史理论最重要的著作《历史的观念》的作者科林伍德（Robin George Collingwood）在该书一开篇，就专门提到该书涉及的历史理论内容不包括对“中国历史思维”的思考。\NoteRef{08}{17}相比西方历史，中国历史记述有两大明显特征：连续的相同文字书写和连续不断的帝系排列。这两点注定了在中国任何时代回看之前的历史，都与当下有连续感。这类似于我们记忆中自己的过去与当下的自己之间的关系。每个人记忆中对自己过往的认知都会对当下自己的行为有影响，而一个人不能将自己的记忆从思想中分离出来，成为思想的客体来进行研究。因此，中国的历史没有被当成过一个客体来进行西方式的研究。


% 原书第356–358页

使用相同文字对历史进行持续记述，首先让中国的历史看起来具有独特的连续感。中国每个朝代的人都能直接读《春秋》，看孔子、庄子的文字。而欧洲各国虽可以从理论上找出与古希腊文明之间的各种联系，但不可否认的是古希腊文和后来古罗马时代使用的拉丁文，以及再往后天主教世界流行的拉丁文都有着根本的区别。文字的不同，使得断裂感在人们打开书的一瞬间就会直观地呈现出来。维特根斯坦（Ludwig Wittgenstein）从哲学层面提出，不同的语言是配合不同的生活形式（form of life）而形成的，一旦形成，则变成这种生活中不可分割的一部分。\NoteRef{08}{18}也就是说，语言以及表达语言的文字是一个族群内维系共有生活形式、互相认同的最基本元素。20世纪以来，学者对美国移民群体的研究，从现实层面揭示了语言对各个移民群体维系原有文化认同所起的关键作用。当移民在历经数代逐渐丧失母语使用能力后，对原有文化的认同也随之丧失。\NoteRef{08}{19}一个读《史记》英文翻译的华裔学生，与直接读《史记》原文的中国学生，尽管二者都能看懂其中历史故事的情节，但在阅读中两人感受到的历史相关性和连续性完全不同。这个不同就是源于语言文字在叙事这个基本功能以外，还附载着文化元素，使阅读原文与阅读译文的人有不一样的主观感受。汉语中有所谓的“感同身受”，字面意思是：尽管未亲身经历，感受就如同亲身经历过一般。但这句话其实有一个逻辑前提，就是要曾经有过类似经历，有过相似体会，才会在看到、听到某种描述时有发生在自己身上一般的感觉。庄子所谓“夏虫不可语冰”，说的就是要有过“感同”，才可能有发生在自己身上一般的感觉。相关性来自感觉。看《史记》原文的中国学生在看到汉字书写历史的第一瞬间就有视觉上的感同；书中项羽感叹“天之亡我，非战之罪”，陈涉的“苟富贵，无相忘”，“燕雀安知鸿鹄之志”，这些话可能早就在电影电视以及生活中听到过，看到时马上就有似曾相识的“同”感；又如“风萧萧兮易水寒，壮士一去兮不复还”，其中苍凉的意境是每一个母语为中文的读者都能心领神会的，并且像这样用环境来塑造气氛的方式也是中文里经常使用，被读者所熟悉和认同的手法。但这样的句子在被翻译后，本来中文要传达出的，以及中文阅读者也能领会到的那层历史意境就会随之减退或者说消失。在英文世界中，这句在一个比较流行的《史记》翻译版本（哥伦比亚大学版）中是这样写的：


% 原书第358页

\begin{bookquote}
Jing Ke came forward and sang this song:\\
Winds cry xiao xiao, Yi waters are cold.\\
Brave men, once gone, Never come back again.\NoteRef{08}{20}
\end{bookquote}

% 原书第358–359页

诗人罗伯特·弗罗斯特（Robert Frost）曾经说：“诗意就是在翻译中丧失的那个部分。（Poetry is what gets left out in translation.）”\NoteRef{08}{21}如果借用这个思路，历史的连续感和相关性也可说是翻译中失去的部分。当英文读者看到“Jing Ke”这个名字，跟我们在中文中看到“恺撒”这个名字时一样，第一反应都是这是“外”人的故事，自然有着不一样的感觉。


% 原书第359页

西方哲学中认为人的认知有两个来源：直觉和逻辑。\NoteRef{08}{22}中国一代又一代使用相同语言文字书写的历史，首先带来的就是直觉上的持续感。同时，中国历史中连续不断的帝王排列又代表了时间上的连续，这就在语言文字的直觉连续以外，又加入了逻辑上的连续。每一个皇帝都是追溯过去时间路径上的标记。作为对比，西方的纪年时间从5世纪后就渐渐统一到了以耶稣诞生年为临界点的公元纪年，因而他们的时间与宗教内部事件相联系，并不依靠世俗君主。“永和九年，岁在癸丑”这是《兰亭集序》的开头。“永和”是晋穆帝司马聃的年号。任何练习书法的人，无论身处哪个时代，都可以通过历史记载来定位自己与王羲之时间上的距离。历史记录对于在中国土地上生活的人，就是一份共有的家谱，可以一直翻到很久以前，并能找到自己与它的联系。同时，中国历史记载也是一份长久的时间表，让人们能够定位出自己与过去某个时间中的人和事之间的距离。


% 原书第359–360页

需要强调的是，这里所说的不仅仅是历史记载上的连续性，也是阅读的人感觉到的连续感。这是主观的感受。这个感受是对过往的一种认同，认为那些人和事与当下有相关性（relevance）。《论语》说：“不患寡而患不均。”讲的是人的感受是在与周围人的比较中而形成的。这是一种横向比较。但有了历史连续感后，中国人的比较维度中还有一个纵向的。李清照表面上看和项羽没有关联，但当金兵南下，她被迫渡江撤离到南方后，感而赋诗：


% 原书第360页

\begin{bookquote}
生当为人杰，死亦为鬼雄。\\
至今思项羽，不肯过江东。
\end{bookquote}

% 原书第360页

李清照从司马迁写的项羽乌江自刎的历史故事中感受到了相关性，并将其视作思考自己人生时的参考。而且，当李清照写下这首诗的时候，她非常清楚项羽之事众所周知，因而她可以用这个典故来表达她诗中想要表达的多重意味。


% 原书第360–361页

“大江东去，浪淘尽，千古风流人物。”历史的连续感，就像词中所言，如江水奔流，未曾断息，直到眼前。苏轼站在赤壁，感受到的是曾经在同一个空间中，那个也曾经掀起风浪的周瑜。几百年前的过往，“遥想公瑾当年，小乔初嫁了”，苏轼谈起来就如同是自己的家事。这时历史的连续感和相关性不仅存在于史书中，同时又通过文学作品进一步被加强。使用历史典故大量存在于中国文学作品中，而在西方文学作品中，典故更多来自《圣经》或者古希腊、古罗马的神话故事。2013年，在“孪生素数猜想”问题取得了重大突破的张益唐，在美国已经生活了二十八年，之前完全不为学界知晓。在突然获得了世界的认可后，接受记者采访时，他引用了一句杜甫的诗“庾信平生最萧瑟，暮年诗赋动江关”，来表达耳顺之年成名的感受。\NoteRef{08}{23}杜甫是在凭吊古迹时，哀叹于自己“漂泊西南天地间”，感怀比他早两百年的庾信的人生际遇。一千多年后的张益唐在诗句中感到了与两个历史人物人生际遇的关联，产生了共鸣。无论身处何方，只要用中文阅读着相同的历史，这种连续感和相关性就会得以维系。语言文字在哪里传递文化，那文化也就在哪里。


% 原书第361–362页

历史叙述并不仅仅影响知识分子，也早已成为大众行为的参考。“孟母三迁”依然会影响到2000年后年轻父母们对住处的选择，这是因为他们感觉到了与孟母的历史故事有相关性，她的做法和思路还有参考价值。再比如本书讲到过的清朝入藏官兵，他们到了拉萨不抢劫财物，而是抢占布达拉宫山崖上的绝佳位置来刻写他们的历史，这就来自以前关于武将的历史叙述对他们的影响。人们在进入一个新的场景时，首先会参看其中跟自己角色类似的人是怎么做的。这是横向的参考。而像清军入藏这种横向维度上缺乏参考的事件，那么历史中的人和事就是纵向的参考。很明显，他们模仿的是“勒石燕然”——效仿东汉窦宪在燕然山大败匈奴后，把胜利的过程刻在石壁上的举动。他们希望自己的故事也能流芳百世。


% 原书第362页

而遗臭万年作为流芳百世的反面，又是所有人避之不及的。作为肯定会留在历史记录中的人，皇帝们一直心怀敬畏的就是如何被后世的历史所书写。他们的一言一行有来自各方面的记录，有点像身边总有各种摄像机跟随，他们知道有些部分的内容会在未来被后世所评论和书写。《春秋》中记载的“齐崔杼弑其君光”，《左传》中以“弑”为中心，写出了崔杼在弑君篡位后，害怕在历史中永留骂名，而逼迫史家篡改历史记录的故事。《左传》中突出的是当时史家不畏死，前仆后继，坚持秉笔直书。但另一方面，这个故事也生动地展示了当权者对在历史中留下永恒污点的担忧。前面章节讲过的皇三子在得知雍正捉拿了皇八子以后，提出他能协助的是：


% 原书第362页

\begin{bookquote}今将杀允禩（注：皇八子）之过落于我身，我情愿承受；主子说我是允禩同党，我也接受。\end{bookquote}

% 原书第362–363页

皇三子明白，已经大权在握的雍正心中唯一的担忧，或者说还不能完全把控的，就是历史评判。《左传》中“郑伯克段于鄢”是当时读书人都熟知的历史。这一段历史批评的是兄弟相残，用“克”字来体现兄弟之情的丧失，郑伯与共叔段兄弟二人，俨然成为了没有关系的路人。《左传》中的这个故事，以及十三经中的历史故事，都是当时人识字启蒙后要反复学习的内容。这些历史故事在皇子以及中国士子中的熟悉程度，绝不亚于《圣经》故事之于西方人。中国的皇帝敬畏在历史中留下骂名，如同西方人担心下地狱。中国历史排列谨严的帝王序列和围绕他们展开的历史书写，让接受了汉文化并熟读中国史书的康熙和他的儿子们对历史的连续感和相关性都有了切实的感受。如何存在于未来永恒的历史书写之中，成为了他们内心深处的忧虑。前文已经讲过，康熙晚年对中国历史上的帝王如数家珍，他排名在什么位置，哪些帝王出过什么问题，他都烂熟于心。皇三子知道雍正担忧即将上演的杀弟桥段将要如何被历史书写。这是独属于他们帝王之家对历史叙述共有的敏感。


% 原书第363页

对于皇帝来说，如何才能保证留在历史上都是正面的记录？要控制历史在未来的叙述，最关键是要让后人接受最初留下的历史资料及最初对这些材料的诠释。那后来的人又如何判断最初的历史叙述是否就是发生过的过去？这里面其实隐含了两个不同的问题：


% 原书第363页

\begin{bookquote}
叙述是否能呈现过去？\\
叙述是否能被相信呈现了过去？
\end{bookquote}

% 原书第363页：原件重复拍摄，正文仅收录一次。

% 原书第364页

前者是一个哲学问题，柏拉图时代就有，而后者是现实操作层面的问题。当陪审团中的所有人做出一致判决的时候，他们就相信了其中一种对过去的叙述。但从哲学层面讲，陪审团相信的那种叙述不见得就真的呈现了过去，只是陪审员们相信了而已。在法庭中，无论是控方还是辩方，他们都是在根据已有的证据，给出一种对过去发生之事的叙述。被陪审团所相信的那种叙述，在法庭上来说，就是还原了过去的叙述。因而皇帝要控制历史的书写，关键点并不在于给出的叙述是否真的还原了过去，而是要保证给出的叙述会被后世相信。


% 原书第364–365页

相信或不相信某条信息或某套叙述，是人的认知中一直在不停处理的问题。按佛教中所说，人的认知首先来自眼、耳、鼻、舌、身、意六根，也就是自己身体直接从外部获取到的信息，然后是间接获取或者说是从学习中得到的认知。这种间接认知有些可以自己验证，比如我们听说南半球的7月是冬天，那么只要在7月时去往南半球就可以验证这一判断。但另一种认知是自己无法验证的知识，比如地球围着太阳转的日心说。我们认可并接受这些自己无法验证的知识，是因为信任了提供该知识的权威。历史作为对过去的叙述，就属于无法亲自验证的一类知识。我们从最初听父母讲祖辈的故事，到后来自己阅读历史，相信某种历史叙述，本质就是在相信该历史的叙述者或者推广这种历史叙述背后的权力机构。新闻报道的对象是刚刚发生的过去。而我们看某条新闻时往往习惯于关注该报道背后的机构，就是在衡量这家机构的信用度，以此决定是否可以相信这则报道。


% 原书第365页

读者相信与否，对于一种历史叙述最后是否被接受至关重要，因此可以看到许多西方史家会专门讲解自己的资料来源，以此来让读者相信自己的叙述是基于可靠的资料。被称为西方史学之父的希罗多德在他的《历史》一书中就反复说明他的信息来源，特别是写到埃及时，他向读者保证所写的内容都是他实地研究的结果，而非基于想象的编造。\NoteRef{08}{24}柏拉图的《理想国》中记载了苏格拉底对说书人（storyteller）和诗人（史诗）的叙事（narrative）的怀疑，认为他们所讲的内容并不可靠，许多都是他们自己的情感表达。\NoteRef{08}{25}过去不能重现，无法像科学中的实验一样，让历史重演一遍。也是这个原因，所有的历史叙述从根本上来讲都有一个是否可信的问题存在。西方历史书写就一直被如何让读者相信这个问题所困扰。本书使用的传教士资料《北京纪事》正是这样，因此，纪理安在书的开头反复申明他的记录都有依据，是发生过的事。


% 原书第365–366页

但在中国传统历史书写中，让读者“相信”并不是一个会经常提及的问题。这源于中国特有的史官传统。大家相信某种历史叙述，源于对官方叙事的信任。中文中有“正史”这个词，也就是暗示这套历史叙述已经被确认，就像“正道”、“正义”一样，不需要怀疑。最早把“正史”当作一个门类的《隋书·经籍志》专门讲述了“正史之缘起”，“古者天子、诸侯”都有国史，而且都有“史官”，这样的历史叙述就是标准，当然就可以直接采信。\NoteRef{08}{26}


% 原书第366页

到了《四库全书》，“正史”成为了史部中的第一类。同时给出的解释是，正史是经过乾隆裁定的史书。这样一来，正史就具有了更明确的官方意味。而且，这时的官方并不是指官方撰写的历史，而是官方认定的历史。像欧阳修的《新五代史》属于私人编撰，只是得到了后世官方的认可，便可跻身“正史”之列。在乾隆确定的二十四史正史中，唐史和五代史都有两套。对于两套并存的问题，《四库全书总目提要》加了一小段话作解释，说这是“钦禀睿裁”的结果。当我们相信正史时，其实就是在相信皇帝的“睿裁”。


% 原书第366页

西方并没有中国这样的史官传统，当然也就没有“正史”的概念，连相对应的词都没有，因而在翻译“正史”一词时，根据“正史”中隐含官方历史的那层意思，翻译为了official history。中国传统中以官方认证的历史为准绳，并不代表传统中国的读者在接受历史叙述时没有产生过怀疑，只是读者已经习惯于把是否“信任”这个本来需要自己主观判断的过程“外包”给了朝廷。


% 原书第366–367页

在这样“信任外包”的传统下，朝廷成为了历史叙述的最终裁定者。特别是到了清朝，连历史编撰中一些关键问题的抉择都由皇帝来作最后定夺。从操作层面来讲，日理万机的皇帝还要亲自来定夺与历史相关的问题，是一件工程浩大的事。但清朝的康雍乾三位皇帝对历史的重视超过了之前的任何朝代。三位皇帝的统治加起来有一百三十八年（1661—1798），超过了存在一百零八年的元朝（1260—1368）。他们三人对历史操控的持续追求，实际上为皇帝争取到了一项以前没有的权力：对历史叙述的决定权。在漫长的中国历史中，皇帝对与自己相关历史的部分进行干预有许多先例，比如玄武门之变后的唐太宗。但他们不会也没有精力过问那些跟自己无关的前朝历史。\NoteRef{08}{27}历史的叙述，对历史的评判，起码在之前的明朝，还掌控在士人手中。但到了清朝，历史评判权逐渐被皇权垄断。


% 原书第367–368页

对历史的全面关注始于康熙。这首先与其个人旺盛的求知欲和操控欲有关。在编纂《明史》的过程中，康熙会亲自阅读审定，多次把编订好的部分发回，命编修们重写；对于编纂中的各种事件专题，他也会亲自过目品评。在平定了三藩之乱的第二年，康熙就组织人写史，用了四年时间纂修而成《平定三逆方略》。后来，在战胜了噶尔丹后，他又组织人撰写了《亲征平定朔漠方略》。这一南一北的两次胜利是康熙最得意的人生功绩，他都在第一时间给出了官方的历史叙述。这样相当于在当朝就推出了事件的正史，定下了对历史事件叙述的基调。从原始材料到历史叙述有一个蜕变的过程。这个过程就是把杂乱无章的原始材料经过筛选与辨识，变成一段有因果逻辑的历史叙述。当一种历史叙事诞生时，往往只有与之相关的原始材料得以保留，而与之不相关的材料就会随着时间被遗忘，然后逐渐消失。像康熙推出了《平定三逆方略》后，三藩方面的许多原始材料就消失了。三藩对抗清廷被称为“三藩之乱”，这一说法沿用至今。后世在沿用这一说法时，其实就已经接受了“乱”字对这一事件的定性。康熙这样第一时间推出“正史”，为事件定调，其实就是在把控历史叙述。后来，这个传统被乾隆发扬光大，他推出了很多部类似的方略。


% 原书第368页

方略这种类型的书都是正面叙述历史功绩，是一种有意识的对以往成绩的总结和推广。而控制历史叙述真正的难点在于如何把负面的部分去除，或者把负面的虚构成正面，同时还能让后世相信。前文已经讲过，康熙十分重视自己在历史中的地位，担心自己由于继承人选择不当，让清朝重蹈隋朝的覆辙，使自己成为历史中的反面教材。康熙就直白地说过他“不畏当时而畏后人”\NoteRef{08}{28}。所谓的“畏后人”就是畏惧后人的历史评判。康熙在指点《明史》撰写的谈话中，也显示出他十分清楚历史叙述的形成，以及后世评判时的各种要素，他说：


% 原书第368–369页

\begin{bookquote}明代实录及纪载事绩诸书，皆当搜罗藏弆。异日《明史》告成之后，新史与诸书俾得并观，以俟天下后世之公论焉。\NoteRef{08}{29}\end{bookquote}

% 原书第369页

前文提到过，按照西方历史理论中的总结，历史叙述的形成受制于三个方面：原始记录、其他版本的历史叙述，以及读者。这里能看到，康熙认为最终的历史叙述，即康熙所谓的“公论”，也是由上面三个方面决定。康熙指挥纂修们收集明代实录（即原始记录）和各种书籍（即其他版本历史叙述），然后收藏好，最后在《明史》编撰完成后，由后世（即读者）对比观看，最后得出公论。\NoteRef{08}{30}在康熙的思维中，很明显，他认为历史叙述的好与坏最终是由后世读者说了算，也就是康熙所说的：


% 原书第369页

\begin{bookquote}《明史》关系极大，必使后人心服乃佳。《宋史》成于元，《元史》成于明，其中是非失实者多，是以至今人心不服。\NoteRef{08}{31}\end{bookquote}

% 原书第369页

康熙说的“后人心服”，就是西方理论中所强调的读者在历史叙述形成中的重要作用。


% 原书第369–370页

读者之所以能发挥作用，前提是假定多种叙述并存，有各种原始资料以供品评。康熙所谓的后世公论，是读者在“新史与诸书”并观下形成。但西方理论和康熙都忽略了另外一种情况，就是当世间只有一种历史叙述和与之匹配的资料留存时读者对历史叙述就失去了影响。


% 原书第370页

雍正一上台就对自己如何被后世书写抱有很强的忧患意识。除了收缴皇十四子与康熙的往来信函，抹去载有皇十四子功绩的石碑，他还命令全部官员上缴康熙时的奏折，不准私藏。这些行为从政治角度可以做出各种解读，但从历史书写角度，这是在控制原始材料流传。而前文讲过的汪景祺和查嗣庭案，本质上就是雍正要抹去其他版本的历史叙述。\NoteRef{08}{32}


% 原书第370页

在没有其他历史叙述的境况下，雍正撰写的历史叙述就会成为唯一。像雍正指控查嗣庭的“日记所载狂妄悖逆之语。与汪景祺相为表里”。除了极少数看过查嗣庭日记的办案人，后世均无从得知其日记究竟有多“大逆不道”。不仅查嗣庭的日记被销毁了，查嗣庭获罪的缘由也无人再敢记录，所有关于这个案子的信息只能从雍正公布的“正规”渠道得知。年羹尧案、皇九子案、皇八子案、隆科多案基本都是如此，现在存世的都只有朝廷给出的历史叙述。


% 原书第370–371页

雍正还曾经提前布局，让整个朝廷在关键时期噤声。没有人说话，或者说评论，那么自然就没有材料流传后世。在汪、查二案之间，雍正逼迫朝廷所有士人声讨曾经参与纂修《明史》的钱名世。钱名世（1660—1730）是康熙四十二年殿试探花，一生基本没有参政，一直在朝廷中从事与编书相关的事务。\NoteRef{08}{33}他善于文词，结交广泛，与京城中大多数名士都有往来。\NoteRef{08}{34}雍正四年（1726）三月，也就是年羹尧被赐死后约三个月，雍正以钱名世曾经在年羹尧鼎盛时期写下多篇诗文赞美年羹尧为由，指责钱名世“称功颂德，备极谄媚”。雍正以“谄媚奸恶，为名教所不容”为名，惩罚钱名世。\NoteRef{08}{35}处罚之一是将钱名世定为“名教罪人”。雍正亲笔写下这四个字，让地方官做成匾额，悬于钱名世在江苏老家的祖宅中。传统中，衣锦还乡、荣归故里是士人心中的向往。雍正则反其道而行之，让钱名世把“名教罪人”的匾额搬回家，辱没家门，使其无颜见江东父老。处罚之二是要求在京官员，只要有举人、进士头衔的，都要作诗一首来谴责钱名世，这便开启了文人互讦的先河。雍正还规定了文体，必须要用诗的形式来批判钱名世。诗言志，诗在传统中是文人用来抒发真实情感的文体，而雍正就是要用写诗这种特别的方式来驯服士人。这种方式不仅有辱斯文，对世道人心的伤害也影响深远。


% 原书第371–372页

钱名世“名教罪人”案发生前两个月，皇八子被抓，雍正宣布了皇八子及其拥护者的各种罪状。\NoteRef{08}{36}雍正四年六月，判定了皇八子“罪状四十款”。\NoteRef{08}{37}钱案发生在决定如何处理皇八子的关键时期。被要求写诗的几百位举人和进士并不属于朝廷的满人权贵圈，也不是参与审判皇八子的核心群体。但这些朝廷中的士子是京城的喉舌，传统上他们是左右世事评判之人。雍正就是要利用“名教所不容”这个借口来驯化群臣，让士人们明白祸从口出、罪由文生的道理；最起码要让他们知道，赞美也可能是一种罪。因为赞美的本质还是对过去的人和事的评价。如果众多像钱名世一样赞美年羹尧的文字都留存了下来，雍正又如何控制对年羹尧的评价？ 2007年在江苏靖江甚至出土了颂扬年羹尧平定青海叛乱的巨大石碑。\NoteRef{08}{38}这个石碑明显是在刻好后，由于政治风向突然变化，才匆忙埋入了地下。简单说来，如果世间存在众多赞美雍正所杀之人的评论，那么雍正又如何将杀了他们的事写入历史？


% 原书第372页

朝廷文臣如此之多，按常理推断，应不乏抗争之人。《永宪录》中提到，“翰林院侍读吴孝登”因所作批判钱名世的诗不合雍正的意思，被判去“宁古塔，给披甲人为奴”；侍读学士陈邦彦被革职。\NoteRef{08}{39}翰林院侍读都是进士出身，是士人中的龙凤。吴孝登能在这样的情况下，逆鳞而上，道德勇气非比寻常，应该也有一段可歌可泣的故事。似乎中国历史上总有吴孝登这样坚决以卵击石、决不妥协的人存在，这也是中国历史读起来总有连续感的又一个原因。吴孝登的故事，后世并没有流传。被发往东北最边远的宁古塔，去给士兵做奴隶，吴孝登曾经的遭遇只能留给后世读史之人凭空想象。桐城派古文大师方苞等都迫于权势写了诗，但这些写过诗的士人都知道这是自己的耻辱，都羞于把这种诗收入自己编订的诗集中。\NoteRef{08}{40}


% 原书第373页

中国历史上“本是同根生，相煎何太急”的案子层出不穷，雍正要消灭他的兄弟，杀他的功臣，放在古今中外政治历史中，也不算特别之事。他要杀年羹尧，要杀皇八子，他就杀了，他有这个权力。但雍正与以前帝王不一样之处在于，他不仅要杀人，还要让后世觉得他杀得合理。所以我们看到这些被杀之人在历史记录中都是经过了公审。他们被杀的罪名都是因为他们罪大恶极，而雍正是在忍无可忍的情况下才杀了他们。年羹尧和皇八子的审判在官方记录中都是“公同议奏”后的结果。雍正不仅要把这些人除去，还要在历史的记录中留下一个公正的美名。


% 原书第373页

雍正无论是合法继位还是篡位称帝，从历史长河中看其实都不稀奇。他的特别之处是他在现实的空间中坐稳皇位后，还非要在历史的空间中再赢一次。他要让后世回看这段历史时，发现他的每一个兄弟都有诸多问题，而他这个皇四子本来就是最应该登位的那一个。康熙说历史叙述的目标是要“后人心服”，雍正就是朝这个目标在努力。他要让所有阅读这段历史的后人心服他，就是康熙心中的继承人。


% 原书第373页

公正地说，雍正对历史的编撰很成功。他的历史叙述在后世被广泛接受。史家陈垣在给钱名世写的小文中，就相信了雍正的逻辑，因为钱名世献媚了年羹尧，才成了“名教罪人”。\NoteRef{08}{41}后世史家也相信皇八子有罪，罪是从康熙朝就积累过来的，也相信了雍正给皇八子扣上的罪名：结党营私。雍正用这个理由来解释为什么有那么多人支持皇八子。尽管也有史家怀疑过雍正对皇八子的指控，但关于皇八子的历史材料和历史叙述既然只有朝廷给出的那一种，那么信也是它，不信也是它。在这种情况下，后世读者没有选择。

% 原书第374页

不过，出乎雍正意料，当时的传教士还写有记录，藏在了欧亚大陆的另一头。这些记录在雍正朝编撰出的历史中划出了一道裂缝，让其中的诸多不实显露了出来。如果没有传教士的记录，那么我们永远不会想到，最早挑起储位之争的人会是雍正本尊，当然也不可能想到与大阿哥最早合谋生事的人也是他。他并不是那个自我标榜的“循理守分，不交结一人，不与闻一事”的皇四子。用传教士的话说，他对“权力的渴望”开启了兄弟相残的夺嫡之争。康熙晚年最后到底发生了什么事可能永远是个谜。“假作真时真亦假，无为有处有还无”，这是曹雪芹对殷鉴未远的家事和国事的感悟。读《大义觉迷录》，看其内容在街头巷尾、远乡僻壤反复宣讲，是曹雪芹、满人权贵以及士子们共同的时代经历。无论书怎么写，改变不了的是亲历者记忆中的过去。除了默默地抹一把辛酸泪，知道真假有无的人又能如何？记忆终将消散，只有文字流传。真事隐、假语存，道出的是一种无奈，但这也是历史存留中不可回避的一种状态。

\begin{sectionnotes}

% 原书注释：原书第439页
\BookNote{08}{1}{节目名称：“Vatican Conclave: Everything You Need to Know about the Meeting to Pick the Next Pope”。}

% 原书注释：原书第439页
\BookNote{08}{2}{“If after three days the cardinals did not arrive at a decision, they were to receive for the next five days only one dish at their noon and evening meals. If these five days elapsed without an election, only bread, wine, and water should be their fare.” Edward A. Pace et al., The Catholic Encyclopedia, v.4, 192.}

% 原书注释：原书第439页
\BookNote{08}{3}{科林伍德（Collingwood）认为历史是在史家提出的问题指导下书写出来，而提出什么问题就是一个主观的选择。R. G. Collingwood, The Idea of History, 9-10；海登·怀特（Hayden White） 认为历史最初来自对原始材料的选择（selection）和整理（arrangement）。Hayden White, Metahistory: The Historical Imagination in Nineteenth- Century Europe, 5.}

% 原书注释：原书第440页
\BookNote{08}{4}{当然这里的“春秋”，也有解释为《春秋》这部书。但后世对这句话的引用中，“春秋”用来泛指历史记录了。}

% 原书注释：原书第440页
\BookNote{08}{5}{《圣祖仁皇帝实录》，卷218。}

% 原书注释：原书第440页
\BookNote{08}{6}{陈寅恪，《柳如是别传》，1250。}

% 原书注释：原书第440页
\BookNote{08}{7}{朱维铮编，《周予同经学史论著选集》，256。}

% 原书注释：原书第440页
\BookNote{08}{8}{Edward Shaughnessy, “ 'New' Evidence on the Zhou Conquest”, 60.孟子提到的《武成》篇就是现存《逸周书》中的《世俘解》，夏含夷该文就是在全面论证《世俘解》是可靠的周初史料。另，最早提出《世俘解》是周初文字的是郭沫若（1930年）。见郭沫若，《中国古代社会研究》，318—320。}

% 原书注释：原书第440页
\BookNote{08}{9}{White, Metahistory: The Historical Imagination in Nineteenth-Century Europe, 5.}

% 原书注释：原书第440页
\BookNote{08}{10}{司马迁与王充在这个问题上的观点，参看Shaughnessy,“New Evidence on the Zhou Conquest”, 60–61。}

% 原书注释：原书第440页
\BookNote{08}{11}{Shaughnessy, “ 'New'Evidence on the Zhou Conquest”, 60.}

% 原书注释：原书第440页
\BookNote{08}{12}{Benedetto Croce, “History Brought Under the General Concept of Art”, 492.}

% 原书注释：原书第440页
\BookNote{08}{13}{Benedetto Croce, Aesthetic as Science of Expression and General Linguistic, 207-209.}

% 原书注释：原书第441页
\BookNote{08}{14}{Benedetto Croce, History as The Story of Liberty, 19–22.}

% 原书注释：原书第441页
\BookNote{08}{15}{北美历史学界对客观性追求，见Peter Novick, That Noble Dream: The "Objectivity Question" and the American Historical Profession.科林伍德在《历史的观念》一书中认为克罗齐理论使得历史学科有了自主性（autonomy），从科学和哲学中分离出来。Collingwood, The Idea of History, 201. 后现代史学的领军人物海登·怀特在采访中承认他的学说受克罗齐影响。Eva Domanska, Encounters: Philosophy of History after Postmodernism, 18.}

% 原书注释：原书第441页
\BookNote{08}{16}{Hayden White, The Content of the Form: Narrative Discourse and Historical Representation, 1987.}

% 原书注释：原书第441页
\BookNote{08}{17}{R. G. Collingwood, The Idea of History, 14.}

% 原书注释：原书第441页
\BookNote{08}{18}{Ludwig Wittgenstein, G. E. M. Anscombe trans., Philosophical Investigations, No.241–243, 88.}

% 原书注释：原书第441页
\BookNote{08}{19}{Joshua Fishman, Language Loyalty in the United States: The Maintenance and Perpetuation of Non-English Mother Tongues by American Ethnic and Religious Groups.}

% 原书注释：原书第441页
\BookNote{08}{20}{Sima Qian, Burton Watson trans. The Records of the Grand Historian: Qin Dynasty, 232. 由 William H. Nienhauser, Jr. 领衔，中美两国学者共同翻译，2021年出版的《史记》英文全译中，这一句的翻译是： "The wind soughs and sighs, The Yi River is cold. Once we hearties leave, We'll never return to our fold!"See, Sima Qian, William H. Nienhauser ed., The Grand Scribe's Records, Revised Volume VII, 616.}

% 原书注释：原书第442页
\BookNote{08}{21}{Peter Robinson, Poetry \& Translation: The Art of the Impossible, 23-24.}

% 原书注释：原书第442页
\BookNote{08}{22}{Croce, Aesthetic as science of expression and general linguistic, 1.}

% 原书注释：原书第442页
\BookNote{08}{23}{\url{https://news.pku.edu.cn/xwzh/52c9644d22314117894466c880452fec.htm}.}

% 原书注释：原书第442页
\BookNote{08}{24}{他在书中经常会说到哪些是他自己的研究，哪些是埃及人或者波斯人告诉他的。比如，他写到埃及部分，专门强调了他写到的内容都是亲自调查后得来，见Herodotus, George Rawlinson trans. History, 162–163。}

% 原书注释：原书第442页
\BookNote{08}{25}{Plato, The Republic，80-82. （中译本见柏拉图著，郭斌和、张竹明译，《理想国》，94—97）}

% 原书注释：原书第442页
\BookNote{08}{26}{《隋书》，卷33。}

% 原书注释：原书第442页
\BookNote{08}{27}{唐宋皇帝已经开始关注历史记录，但他们都还是主要关注跟他们自己相关事情的记录。见蔡崇榜，《宋代修史制度研究》，193—195。}

% 原书注释：原书第442页
\BookNote{08}{28}{《圣祖仁皇帝实录》，卷218。}

% 原书注释：原书第442页
\BookNote{08}{29}{《圣祖仁皇帝实录》，卷154。}

% 原书注释：原书第442页
\BookNote{08}{30}{在康熙思路下撰写的《明史》，不再把问题归结到皇帝那里，而是归咎于大臣身上，加入了《奸臣传》这个门类。对于这样写史，当代的明史专家认为有失公允。参看顾诚，《论钦定〈明史〉作〈奸臣传〉之微意—兼为严嵩等人重新定位》。}

% 原书注释：原书第442页
\BookNote{08}{31}{《圣祖仁皇帝实录》，卷218。}

% 原书注释：原书第442页
\BookNote{08}{32}{以前说到这些案子时，都算作“文字狱”的一种。对“文字狱”的传统理解是清朝为了防止反清复明，着眼于满人巩固统治这个角度。这样的解释忽略了许多被归结为“文字狱”案子完全没有反对满人政权的意味。像汪景祺和查嗣庭都是真心归附清朝，而且希望在清朝廷干出一番事业的人。他们都拥护清朝政权。他俩获罪完全是因为皇位更替后，皇帝对私自记录历史的标准改变了，他们才获得大罪。二人都是有高度政治敏感的老手，他们当时写下的东西都是他们认为可以写下的。问题是他们写下时默认用到的是康熙朝的标准，而雍正改变了这个标准，他们就成为了被处理的典型。无论从政治角度，还是文字角度入手，以前研究都无法解释为什么汪、查两个不相关的人会被雍正算作一类，反复把他俩联系到一起。只有从历史撰写角度，也就是二人记录下了雍正认为不应该记录下的历史入手，才能理解雍正对二人严酷处理的缘由。雍正用公开的审理、祸及族人的方式，来杜绝士人私下记载朝中之事的行为。雍正要从源头上制止其他版本的历史叙述出现。雍正对二人的处罚也一致：枭首示众，以儆效尤。}

% 原书注释：原书第443页
\BookNote{08}{33}{钱名世仕途，参见余子铨，《“名教罪人”案与常州钱氏家族》，47—55。钱名世在康熙朝还写了一部《四藩本末》（三藩之外加入了孔有德）。钱名世，《四藩本末》，《明清史料丛书续编》，第2册。}

% 原书注释：原书第443页
\BookNote{08}{34}{钱名世交友，参见余子铨，《“名教罪人”案与常州钱氏家族》，36—47。}

% 原书注释：原书第443页
\BookNote{08}{35}{雍正四年三月壬戌（1726年5月1日），《世宗宪皇帝实录》，卷42。}

% 原书注释：原书第444页
\BookNote{08}{36}{“并令宗人府、将允禩、允禟、苏努、吴尔占名字除去”，《世宗宪皇帝实录》，卷40。}

% 原书注释：原书第444页
\BookNote{08}{37}{“康亲王崇安，及诸王、贝勒、贝子、公、满汉文武大臣等，公同议奏。阿其那罪状四十款”，《世宗宪皇帝实录》，卷45。}

% 原书注释：原书第444页
\BookNote{08}{38}{常永明，《试解靖江孔庙巨碑之谜》。}

% 原书注释：原书第444页
\BookNote{08}{39}{《永宪录》，卷4。}

% 原书注释：原书第444页
\BookNote{08}{40}{陈垣，《钱名世轶事》，85。}

% 原书注释：原书第444页
\BookNote{08}{41}{陈垣，《钱名世轶事》。}

\end{sectionnotes}
